\hypertarget{low-level-networking-and-sockets}{%
\chapter{Low-Level Networking and
Sockets}\label{low-level-networking-and-sockets}}

Networking in Python is built upon the foundational Berkeley Sockets
API. While high-level libraries like \passthrough{\lstinline!requests!}
or \passthrough{\lstinline!httpx!} are common, systems engineering
requires mastery of the low-level \passthrough{\lstinline!socket!} and
\passthrough{\lstinline!ssl!} modules.

\hypertarget{socket-the-berkeley-interface}{%
\subsection{\texorpdfstring{38.1 \texttt{socket}: The Berkeley
Interface}{38.1 socket: The Berkeley Interface}}\label{socket-the-berkeley-interface}}

A socket is an endpoint for communication. The
\passthrough{\lstinline!socket!} module provides a C-like interface to
the operating system's networking stack.

\hypertarget{address-families-and-socket-types}{%
\subsubsection{1. Address Families and Socket
Types}\label{address-families-and-socket-types}}

\begin{itemize}
\tightlist
\item
  \textbf{AF\_INET / AF\_INET6}: IPv4 and IPv6 networking.
\item
  \textbf{AF\_UNIX}: Unix Domain Sockets (local IPC, faster than network
  sockets as they skip the TCP/IP stack).
\item
  \textbf{SOCK\_STREAM}: TCP (reliable, connection-oriented).
\item
  \textbf{SOCK\_DGRAM}: UDP (unreliable, connectionless).
\end{itemize}

\hypertarget{the-lifecycle-of-a-server-socket}{%
\subsubsection{2. The Lifecycle of a Server
Socket}\label{the-lifecycle-of-a-server-socket}}

\begin{enumerate}
\def\labelenumi{\arabic{enumi}.}
\tightlist
\item
  \textbf{\passthrough{\lstinline!socket()!}}: Create the socket
  descriptor.
\item
  \textbf{\passthrough{\lstinline!bind()!}}: Associate the socket with
  an address and port.
\item
  \textbf{\passthrough{\lstinline!listen()!}}: Enable the socket to
  accept connections (sets the backlog size).
\item
  \textbf{\passthrough{\lstinline!accept()!}}: Block until a client
  connects. Returns a \textbf{new} socket object specifically for that
  connection.
\end{enumerate}

\hypertarget{blocking-vs.-non-blocking}{%
\subsubsection{3. Blocking
vs.~Non-blocking}\label{blocking-vs.-non-blocking}}

By default, sockets are blocking. Setting
\passthrough{\lstinline!sock.setblocking(False)!} makes
\passthrough{\lstinline!send!} and \passthrough{\lstinline!recv!} return
immediately, raising \passthrough{\lstinline!BlockingIOError!} if no
data is available. This is the foundation for multiplexing (as seen in
Chapter 10).

\hypertarget{ssl-secure-communication}{%
\subsection{\texorpdfstring{38.2 \texttt{ssl}: Secure
Communication}{38.2 ssl: Secure Communication}}\label{ssl-secure-communication}}

The \passthrough{\lstinline!ssl!} module wraps OpenSSL to provide
TLS/SSL encryption.

\hypertarget{sslcontext}{%
\subsubsection{\texorpdfstring{1.
\texttt{SSLContext}}{1. SSLContext}}\label{sslcontext}}

This object stores configuration (certificates, cipher suites, protocol
versions). * \textbf{Certificate Verification}:
\passthrough{\lstinline!context.verify\_mode!} ensures the server's
identity is valid against a CA bundle. * \textbf{ALPN/SNI}: Support for
modern TLS features like Application-Layer Protocol Negotiation (used
for HTTP/2) and Server Name Indication.

\hypertarget{wrapping-sockets}{%
\subsubsection{2. Wrapping Sockets}\label{wrapping-sockets}}

You don't create an SSL socket directly; you ``wrap'' an existing TCP
socket:

\begin{lstlisting}[language=Python]
conn = context.wrap_socket(raw_sock, server_hostname="example.com")
\end{lstlisting}

This triggers the TLS handshake process.

\hypertarget{mmap-memory-mapped-files}{%
\subsection{\texorpdfstring{38.3 \texttt{mmap}: Memory-Mapped
Files}{38.3 mmap: Memory-Mapped Files}}\label{mmap-memory-mapped-files}}

\passthrough{\lstinline!mmap!} allows you to map a file directly into
the process's virtual memory space.

\hypertarget{why-use-mmap}{%
\subsubsection{\texorpdfstring{1. Why use
\texttt{mmap}?}{1. Why use mmap?}}\label{why-use-mmap}}

\begin{itemize}
\tightlist
\item
  \textbf{Performance}: Reading from an \passthrough{\lstinline!mmap!}
  object is often faster than standard \passthrough{\lstinline!read()!}
  calls because it avoids copying data from kernel space to user space
  (Zero-copy).
\item
  \textbf{IPC}: Multiple processes can map the same file. Changes made
  by one process are immediately visible to others, providing a
  high-speed shared memory mechanism.
\end{itemize}

\hypertarget{interface}{%
\subsubsection{2. Interface}\label{interface}}

\passthrough{\lstinline!mmap!} objects behave like both a bytearray and
a file. They support slicing, regex searching, and standard
\passthrough{\lstinline!read!}/\passthrough{\lstinline!write!} methods.

\hypertarget{performance-optimizations-sendfile}{%
\subsection{\texorpdfstring{38.4 Performance Optimizations:
\texttt{sendfile}}{38.4 Performance Optimizations: sendfile}}\label{performance-optimizations-sendfile}}

For high-performance file serving, Python provides
\passthrough{\lstinline!os.sendfile!}. * \textbf{Zero-Copy}: It
instructs the kernel to copy data directly from a file descriptor (disk)
to a socket descriptor (network) without the data ever entering the
Python interpreter's memory. This drastically reduces CPU usage and
memory bandwidth for static file delivery.

\begin{center}\rule{0.5\linewidth}{0.5pt}\end{center}

\begin{center}\rule{0.5\linewidth}{0.5pt}\end{center}
